%%%%%%%%%%%%%%%%%%%%%%%%%%%%%%%%%%%%%%%%%%%%%%%%%
\section{Training-Time Data Augmentation}\label{sec:data_augmentation}

\Cref{sec:data_sourcing_and_taxonomy,sec:data_quality_and_curation,sec:synthetic_data_supplementation} describe how the dataset was constructed. This section describes how it is transformed at training time. Augmentation is treated as a variable to control for in service of the central question, whether distillation outperforms direct fine-tuning, rather than as an object of study in its own right.

%%%%%%%%%%%%%%%%%%%%%%%%%%%%%%%%%%%%%%%%%%%%%%%%%
\subsection{Augmentation as a Controlled Variable}\label{sec:augmentation_controlled_variable}

Two properties of an augmentation transform determine how it may be used here. The first is \emph{distillation compatibility}. Single-image transforms such as flipping, color jitter and affine scaling preserve a one-to-one correspondence between an image and its augmented version, so the same augmented image can be shown to both teacher and student. Multi-image compositing transforms such as mosaic, mixup and copy-paste build a new image from fragments of several sources, and no well-defined single teacher view exists for such a composite: the teacher would have to reconcile predictions for several unrelated crops in one frame. Compositing is therefore never used in the distillation condition. The second property is \emph{architecture specificity}. Flipping, color jitter, scale and translation are architecture-agnostic, whereas mosaic is tuned to YOLO's anchor-coverage assumptions and does not transfer unchanged to an arbitrary detector.

%%%%%%%%%%%%%%%%%%%%%%%%%%%%%%%%%%%%%%%%%%%%%%%%%
\subsection{Experimental Setups}\label{sec:augmentation_setups}

Three training configurations share a fixed student architecture (\textit{YOLO26n} \cite{UltralyticsYOLO262025}) and differ only in loss function and augmentation policy, as summarized in \Cref{tab:augmentation-setups}.

\begin{table}[H]
\centering
\caption{Experimental setups sharing a fixed student architecture.}
\label{tab:augmentation-setups}
\begin{tabularx}{\textwidth}{@{}clXX@{}}
\toprule
Setup & Loss & Augmentation & Role \\
\midrule
A & Teacher$\rightarrow$student distillation & Basic set only (\Cref{sec:basic_augmentation}) & Distillation condition \\
B & Direct detection & Basic set only, byte-identical to Setup A & Direct-fine-tuning baseline \\
C & Direct detection & Basic set $+$ compositing (\Cref{sec:compositing_augmentation}) & Student's augmentation ceiling \\
\bottomrule
\end{tabularx}
\end{table}

Setup A against Setup B isolates the benefit attributable to distillation itself, since the two are identical except in the training signal. This is the headline comparison the work is built around, and it requires the two augmentation pipelines to be byte-identical rather than nominally the same. Setup B against Setup C isolates the benefit of compositing on top of direct training. Setup A against Setup C answers a practitioner's question, whether distillation with modest augmentation beats direct training at the student's full augmentation budget, and it is not a clean single-factor comparison, since it varies the training signal and the augmentation policy at once.

%%%%%%%%%%%%%%%%%%%%%%%%%%%%%%%%%%%%%%%%%%%%%%%%%
\subsection{Basic Augmentation Set (Setups A and B)}\label{sec:basic_augmentation}

\Cref{tab:basic-augmentation} lists the transforms in the basic set shared by Setups A and B, and those considered and excluded.

\begin{table}[H]
\centering
\caption{Basic augmentation set used identically in Setups A and B.}
\label{tab:basic-augmentation}
\begin{tabularx}{\textwidth}{@{}lcX@{}}
\toprule
Transform & Included & Rationale \\
\midrule
Horizontal flip ($p=0.5$) & Yes & Label-preserving, doubles effective viewpoint coverage \\
HSV color jitter (H $0.015$, S $0.7$, V $0.4$) & Yes & Robustness to lighting and color variation, no geometric change \\
Random scale ($\pm 0.5$), translation ($0.1$) & Yes & Taken from the tested \textit{YOLOv5} \cite{ultralyticsComprehensiveGuideUltralytics} \texttt{random\_perspective} routine \\
Rotation, shear, perspective warp & No & Risk of bounding-box distortion outweighs the expected benefit \\
Vertical flip & No & Terrestrial mammals are consistently upright \\
Blur, CutOut, grayscale & No & Kept out to hold the basic set minimal \\
\bottomrule
\end{tabularx}
\end{table}

%%%%%%%%%%%%%%%%%%%%%%%%%%%%%%%%%%%%%%%%%%%%%%%%%
\subsection{Compositing Augmentation (Setup C)}\label{sec:compositing_augmentation}

Setup C adds three compositing transforms on top of the basic set. Mosaic tiles four training images into one frame at probability $\approx 1.0$, following \textit{YOLOv4} convention \cite{bochkovskiyYOLOv4OptimalSpeed2020}. Mixup linearly blends two images and their labels at probability $\approx 0.1$ \cite{zhangMixupEmpiricalRisk2018}. Copy-paste \cite{ghiasiSimpleCopyPasteStrong2021} is implemented but is a no-op here, since the curated ground truth carries bounding boxes only and no masks to paste. Mosaic and mixup are disabled for the final $10\%$ of epochs, the \enquote{close-mosaic} tail, so the model fine-tunes on the un-composited distribution before training ends.

Compositing is not a substitute for real training data. Augmenting a small image set many times over fits a model to those images' poses and backgrounds rather than to the underlying class, which is the reasoning behind the $100$-real-image floor in the Band C control condition of \Cref{sec:band_structure}. Augmentation multiplies the variance of an adequate pool. It does not stand in for one.

%%%%%%%%%%%%%%%%%%%%%%%%%%%%%%%%%%%%%%%%%%%%%%%%%
\subsection{Reproducibility and Confound Control}\label{sec:augmentation_reproducibility}

Because the Setup A/B comparison is the work's central result, the two configurations share a single dataset and transforms code path rather than parallel implementations, so any difference in outcome is attributable to the teacher rather than to a discrepancy between two pipelines. Seeding is controlled so that the same augmented image reaches teacher and student within a given distillation step, and an assertion verifies that the compositing flags are disabled before a distillation run begins. Every run's augmentation configuration is logged alongside its training metrics (\Cref{sec:model_training_experiment_tracking}), and Band A test mAP is logged separately, since classes trained entirely on synthetic imagery are the most likely to reveal augmentation amplifying generation artifacts.
